Models that exceed GPU memory can be served by keeping some layers resident and streaming the others from host-attached memory, such as page-locked DRAM or a CXL expander, through a ring of GPU buffers. Where should the resident layers sit? We model the pipeline as a max-plus linear system whose iteration time is the maximum mean of a family of tilings of the layer sequence. For $L$ layers of equal compute time $c$ and copy time $f$, $r$ resident, $m$ streamed and $k\le m$ buffers we prove three results: resident layers are useless when a layer computes at least as long as it copies; otherwise, spreading the resident layers evenly attains the lower bound $\max(Lc,mf)$ whenever $(k-1)f\ge 2c$, so three buffers suffice; and a contiguous block of resident layers, as produced by model-order filling, takes $\max(Lc,\,mf,\,rc+2c+(m-k)f)$, up to $2-c/f$ times the optimum. On a Tesla T4 with page-locked copies at \nHwBw~GB/s, similar to current CXL expanders, real Llama-3.2-3B and Llama-3.1-8B layer kernels and per-layer times measured on that GPU, the model reproduces \nAccN{} measured iteration times with a mean error of \nAccMape\% and orders \nAccRankExact\% of distinguishable placement pairs correctly. Even spreading is up to \nMeasPenaltyMax$\times$ faster than a contiguous block and reaches the bound with two buffers where the block needs five. Because the T4 is power-limited, the main measurements lock its clock.
